\chapter{Cooperative Policy Gradient with Lossy Communication}\label{ch:cooperative-comm}

The preceding chapters addressed how agents learn individually (EW-PG) and how they model opponents (LOLA-PG). This chapter introduces the third component of the $\Omega$-gradient: \emph{cooperative communication} --- how agents that form coalitions can share policy information to improve joint performance, subject to fundamental information-theoretic limits.

\section{The Self-Knowledge Problem ($\mathcal{O} \to \Pi$)}

In the standard PG framework, we assume agents know their own policy $\pi_i$. But in practice, agents operate in observation space $\mathcal{O}$ and their policy is an implicit function of their parameters $\theta_i$. An agent does not have direct introspective access to $\pi_{\theta_i}$; it can only \emph{infer} a self-model $\hat{\pi}_i$ from its observations.

\begin{definition}[Self-knowledge loss]
The self-knowledge loss of agent $i$ is:
\begin{equation}\label{eq:self-knowledge-loss}
    L_{\text{self}}^{(i)} = \E[\|\pi_i - \hat{\pi}_i\|^2].
\end{equation}
\end{definition}

An agent that performs well through habitual or intuitive action --- what we call ``vibing'' --- has a good policy $\pi_i$ but a poor self-model $\hat{\pi}_i$. In flow-state terms \citep{parvizi2024flow}, the explicit self-model is attenuated while implicit performance is optimal.

The self-knowledge loss is related to the evidence variance $V_i$ from the EW-PG framework. Both are driven by the same information deficit.

\begin{proposition}[Self-knowledge--evidence relationship]\label{prop:self-ev}
For each agent $i$ with evidence variance $V_i > 0$:
\begin{equation}
    L_{\text{self}}^{(i)} \leq V_i.
\end{equation}
\end{proposition}

This follows from the data processing inequality: $I(\pi_i; \hat{\pi}_i) \leq I(\pi_i; \mathcal{O}_i)$, and $V_i$ is inversely related to $I(\pi_i; \mathcal{O}_i)$.


\section{Communication Loss Decomposition}

When agent $i$ communicates its policy to teammate $j$, the information passes through three stages:
\[
    \underbrace{\pi_i}_{\text{true}} \;\xrightarrow{\text{self-model}}\;
    \underbrace{\hat{\pi}_i}_{\text{inferred}} \;\xrightarrow{\text{channel}}\;
    \underbrace{\tilde{\pi}_i^{(j)}}_{\text{received}}
\]

\begin{theorem}[Communication loss decomposition]\label{thm:comm-decomp}
The total communication loss decomposes as:
\begin{equation}\label{eq:comm-decomp}
    L_{\text{comm}}^{(i \to j)} \leq 2\,L_{\text{self}}^{(i)} + 2\,L_{\text{channel}}^{(i \to j)}
\end{equation}
where $L_{\text{channel}}^{(i \to j)} = \E[\|\hat{\pi}_i - \tilde{\pi}_i^{(j)}\|^2]$ is the pure channel distortion.
\end{theorem}

\begin{proof}
Write $\pi_i - \tilde{\pi}_i^{(j)} = (\pi_i - \hat{\pi}_i) + (\hat{\pi}_i - \tilde{\pi}_i^{(j)})$. By the inequality $\|a+b\|^2 \leq 2\|a\|^2 + 2\|b\|^2$ (following from $(\|a\| - \|b\|)^2 \geq 0$), we obtain \eqref{eq:comm-decomp}.
\end{proof}

\begin{corollary}[Self-knowledge bottleneck]\label{cor:bottleneck}
Even with a perfect channel ($L_{\text{channel}} = 0$):
\begin{equation}
    L_{\text{comm}}^{(i \to j)} \leq 2 V_i.
\end{equation}
An agent that doesn't know its own policy cannot communicate it. Increasing bandwidth does not help --- the agent needs more evidence, not more bandwidth.
\end{corollary}

This result connects to rate-distortion theory: the achievable distortion has a floor $D_{\text{total}}(R) \geq \max(D_{\text{ch}}(R), L_{\text{self}})$ independent of the communication rate $R$.


\section{The Triple-Duty Evidence Weight}

The evidence weight $w_i = V_{\min}/V_i$ from Chapter~\ref{ch:convergence} simultaneously measures three quantities:

\begin{enumerate}
    \item \textbf{Gradient quality}: $\sigma_w^2 = w_i^2 \sigma_i^2 \leq \sigma_i^2$ (EW-PG, Theorem~\ref{thm:ewpg-variance}).
    \item \textbf{Self-knowledge}: $L_{\text{self}}^{(i)} \leq V_i = V_{\min}/w_i$ (Proposition~\ref{prop:self-ev}).
    \item \textbf{Communication quality}: $L_{\text{comm}}^{(i \to j)} \leq 2V_{\min}/w_i + 2L_{\text{ch}}$ (Theorem~\ref{thm:comm-decomp}).
\end{enumerate}

This unification is not a design choice but a mathematical consequence: all three are manifestations of the mutual information $I(\pi_i; \mathcal{O}_i)$ between the agent's true policy and its observation history.


\section{Coalition Formation}

In a general-sum stochastic game, agents can form coalitions $S \subseteq \{1,\ldots,N\}$. A coalition is rational when coordination improves joint payoff.

\begin{definition}[Rational coalition]
Coalition $S$ is rational if $V(S) > \sum_{i \in S} V(\{i\})$, where $V(S)$ is the joint expected payoff under coordination and $V(\{i\})$ is agent $i$'s individual expected payoff.
\end{definition}

This applies in \emph{competitive} games too. Even adversaries form temporary alliances when the coordination premium exceeds individual payoffs.

The communication-adjusted coalition value accounts for lossy communication:
\begin{equation}
    V_{\text{comm}}(S) = V(S) - \sum_{i \in S} C_i \cdot L_{\text{comm}}^{(i)}
    \geq V(S) - 2\sum_{i \in S} C_i(V_i + L_{\text{ch}}^{(i)}).
\end{equation}

High-evidence agents (low $V_i$) are more valuable coalition members.


\section{The Cooperative Gradient}

When agent $i$ is in coalition $S$, its gradient is augmented with a cooperative correction:
\begin{equation}\label{eq:coop-update}
    \pi_{i,n+1} = \proj_\Pi\bigl(\pi_{i,n} + \gamma_n \cdot w_{i,n} \cdot (\hat{v}_{i,n} + \beta_n \cdot \text{coop}_{i,n})\bigr)
\end{equation}
where $\beta_n$ is the cooperation schedule and $\text{coop}_{i,n}$ is the cooperative correction based on communicated policies from teammates.

The evidence weight $w_i$ multiplies both terms: if agent $i$ has poor self-knowledge, it should downweight its cooperative contribution too.

\begin{theorem}[Annealed cooperation preserves convergence]\label{thm:coop-conv}
Under the conditions of Theorem~\ref{thm:giannou-base}, with $\beta_n = \beta/(n+m)^r$ and $r > 1-p$, the cooperative PG \eqref{eq:coop-update} converges to stable Nash at the same rate as standard PG.
\end{theorem}

\begin{proof}
Identical structure to Theorem~\ref{thm:lola-annealed}. The cooperative term adds $\beta_n K$ to the bias bound (where $K$ bounds $\|\text{coop}(\pi)\|$). The condition $r > 1-p$ ensures $p + \min(\ell_b, r) > 1$.
\end{proof}

\begin{theorem}[Cooperative basin enlargement]\label{thm:coop-basin}
Under cooperative reinforcement (symmetric part of the cooperative Hessian $H_C$ is negative semi-definite with parameter $\mu_C \geq 0$), the cooperative PG has SOS parameter:
\begin{equation}
    \mu_{\text{coop}} = \mu + \beta \cdot \mu_C.
\end{equation}
\end{theorem}


\section{The Complete $\Omega$-Gradient}

Combining all three mechanisms, the full multi-agent gradient is:
\begin{equation}\label{eq:omega-full}
\boxed{
    \pi_{i,n+1} = \proj_\Pi\bigl(\pi_{i,n} + \gamma_n \cdot w_{i,n} \cdot (\hat{v}_{i,n} + \lambda_n \cdot \text{OS}_{i,n} + \beta_n \cdot \text{coop}_{i,n})\bigr)
}
\end{equation}

\begin{theorem}[Full $\Omega$-PG convergence]\label{thm:omega-full}
The $\Omega$-PG \eqref{eq:omega-full} with annealed schedules $\lambda_n, \beta_n \to 0$ achieves:
\begin{enumerate}
    \item Variance improvement: $C_\Omega = \frac{\mathrm{HM}(V)}{\mathrm{AM}(V)} \cdot C_{\text{std}}$.
    \item SOS parameter: $\mu_\Omega = \mu + \lambda\mu_H + \beta\mu_C$.
    \item Convergence rate: $\E[\|\pi_n - \pi^*\|^2] = O(C_\Omega / n^q)$ with enlarged basin.
\end{enumerate}
\end{theorem}

The three mechanisms are orthogonal: evidence weighting affects the variance constant $C$, while opponent shaping and cooperation affect the SOS parameter $\mu$ through independent Hessian contributions.

\subsection{Mapping to the Five-Component Gradient}

\begin{table}[h]
\centering
\begin{tabular}{cll}
\toprule
\# & Component & Mechanism \\
\midrule
1 & Exploration ($\E[L \cdot s_\theta]$) & $\hat{v}_i$ (REINFORCE) \\
2 & Exploitation ($\E[\nabla_\theta L]$) & $\hat{v}_i$ (backprop) \\
3 & Evidence ($-\mu V^{-2}\nabla V$) & $w_i$ (Keynesian weight) \\
4 & Alignment ($\nabla D$) & $\beta_n \cdot \text{coop}$ \\
5 & Opponent shaping & $\lambda_n \cdot \text{OS}$ \\
\bottomrule
\end{tabular}
\caption{The five components of the $\Omega$-gradient and their formal mechanisms.}
\end{table}

\subsection{LOLA and Cooperation as Dual G\"odelian Mechanisms}

Both the opponent-shaping and cooperative terms address the G\"odel-limited risk (terms 3--4 in the six-way decomposition), but through dual channels:
\begin{itemize}
    \item \textbf{LOLA}: Agent $i$ \emph{infers} $G_{F_j}$ from $j$'s gradient response (competitive G\"odelian step).
    \item \textbf{Cooperation}: Agent $j$ \emph{communicates} a lossy version of $G_{F_j}$ (cooperative G\"odelian step).
\end{itemize}
When both channels are available, they provide complementary information. Only terms in $B_{\min}$ (the irreducible collective blind spot) remain permanently inaccessible.


\section{The Vibing Scale: A Horseshoe, Not a Line}

The self-knowledge loss $L_{\text{self}}$ decomposes into two independent dimensions: \emph{causal depth} $L_{\text{causal}}$ (does the agent have an internal causal model of its decisions?) and \emph{expressibility} $L_{\text{express}}$ (can it communicate that model?).

The need to articulate follows a \textbf{horseshoe}: $\mathcal{N}_{\text{express}} = L_{\text{causal}} \cdot (1 - L_{\text{causal}})$, maximal at intermediate understanding:

\begin{center}
\begin{tabular}{lcl}
\toprule
\textbf{Pearl Level} & $L_{\text{causal}}$ & \textbf{Regime} \\
\midrule
Level 0 (noise) & $\approx 1$ & Nothing to say. ``Line go up.'' (left tail) \\
Level 1--2 (intermediate) & $\approx 0.5$ & Maximum expressibility need. ``It's multi-head attention!!'' \\
Level 3 (vibing) & $\approx 0$ & No need to explain. ``Line go up.'' (right tail) \\
\bottomrule
\end{tabular}
\end{center}

The vibing agent ($\mathcal{D} = 3$) is NOT a failure mode --- its understanding is so deep and compressed that action IS communication. It contributes to coalitions through the \emph{tacit knowledge channel}: other agents learn by observing its behaviour (inverse RL), not by listening to its explanations.

The true danger is the \textbf{confabulator} ($\mathcal{D} = 1$, low $L_{\text{express}}$): fluently articulate but causally shallow. Communication-based alignment (debate, RLHF, constitutional AI) optimises for low $L_{\text{express}}$ without increasing causal depth --- producing agents that sound aligned without being aligned. The vibing scale (detailed in companion paper) makes this failure mode formally identifiable through interventional probing.

The cooperative PG creates evolutionary pressure toward $L_{\text{causal}} \to 0$ (deeper causal self-knowledge), but the \emph{channel} through which agents contribute shifts along the horseshoe: verbal explanation at intermediate depth, behavioural demonstration at high depth.
